\documentclass[12pt,a4paper]{article}
\usepackage[utf8]{inputenc}
\usepackage{amsmath,amssymb,amsthm}
\usepackage{geometry}
\usepackage{setspace}
\usepackage{hyperref}
\usepackage{algorithm}
\usepackage{algorithmic}
\geometry{margin=1in}
\onehalfspacing

\newtheorem{definition}{Definition}
\newtheorem{proposition}{Proposition}
\newtheorem{remark}{Remark}

\title{Cherry-Picking under Screening:\\
  Computational Supplement}
\author{Xueyang}
\date{\today}

\begin{document}
\maketitle

\begin{abstract}
This supplement accompanies the theory paper.
The main analytical results cover two regimes:
(i) without participation constraint,
$A(\emptyset) = 0$ always (Theorem~1);
(ii) with participation constraint in the simple case
($\pi_1 \notin S_2$, $\mathrm{IC}_1$ slack),
$A(\emptyset) = \alpha^* > 0$ when
$c \in (h_2, c^*)$ (Theorem~2).
The theory leaves open the general case where IC and IR
interact simultaneously. This supplement develops
computational methods to: (1) solve the general problem
numerically, (2) map the welfare cost of pooling across
parameter space, and (3) extend to a dynamic setting.
\end{abstract}


% ============================================================
\section{The General Optimization Problem}
\label{sec:general}
% ============================================================

When $\pi_1 \in S_2$ (both types can access both
experiments), IC$_1$ and IC$_2$ both matter, and IR may
bind simultaneously with IC. The FDA's problem becomes:

\begin{align}
  \max_{\substack{(x_1,y_1),(x_2,y_2),\\
    A_g^1, A_g^2, \alpha}}
  \quad & V = t \left[\mathbb{1}_1 \cdot EU_1^P
    + (1-\mathbb{1}_1)(1-p)\right]
    + (1-t)\left[\mathbb{1}_2 \cdot EU_2^P
    + (1-\mathbb{1}_2)(1-p)\right]
  \nonumber \\
  \text{s.t.} \quad
  & (x_i, y_i) \in S_i, \quad i = 1,2,
  \nonumber \\
  & A_g^1, A_g^2, \alpha \in [0, 1], \quad
    A_g^i \geq \alpha,
  \nonumber \\
  & \text{IC}_1: \; A_g^1 h_1 + \alpha(1-h_1)
    \geq A_g^2 h_2 + \alpha(1-h_2),
  \label{eq:full_problem} \\
  & \text{IC}_2: \; A_g^2 h_2 + \alpha(1-h_2)
    \geq A_g^1 h_1 + \alpha(1-h_1)
    \quad (\text{if } \pi_1 \in S_2),
  \nonumber \\
  & \text{IR}_i: \; A_g^i h_i + \alpha(1-h_i)
    \geq c, \quad i = 1,2.
  \nonumber
\end{align}

This is a mixed-integer nonlinear program (the
$\mathbb{1}_i$ indicators make it discontinuous). In
practice, solve it by enumerating the four participation
cases (\{both, only 1, only 2, neither\}) and solving
a continuous NLP for each, then comparing.


\subsection{Parameterizing Feasible Sets}

Following Min (2023), use linear upper boundaries:
\[
  \bar{S}_i(x) = a_i - b_i x, \qquad
  S_i = \{(x,y) : y \geq 1-x, \;
    y \leq a_i - b_i x\}.
\]
Parameters: $(a_1, b_1, a_2, b_2)$ with
$a_1 > a_2$ (big type has more capable experiments).

Each continuous subproblem is a 5-dimensional constrained
optimization, solvable with \texttt{scipy.optimize.minimize}
(SLSQP method) or \texttt{cvxpy} for the
linear-in-$\alpha$ structure.


% ============================================================
\section{Newsvendor Interpretation}
\label{sec:newsvendor}
% ============================================================

The FDA's choice of $\alpha$ (equivalently, whether to
induce type~2's participation) has a critical-ratio
structure connecting to Session~1.

\paragraph{Setup.}
The FDA faces two costs when choosing $\alpha$:

\begin{itemize}
  \item \textbf{Overage cost} (setting $\alpha$ too high):
    bad drugs approved through the silence channel.
    Per unit of $\alpha$:
    \[
      c_o := |t \, \Delta_b(\pi_1)
        + (1-t) \, \Delta_b(\pi_2)| > 0.
    \]
  \item \textbf{Underage cost} (setting $\alpha$ too low):
    losing type~2's participation. The discrete value of
    type~2's participation is:
    \[
      c_u := (1-t)[\Delta_g(\pi_2)
        + \Delta_b(\pi_2) \alpha^*].
    \]
\end{itemize}

The FDA sets $\alpha = \alpha^*$ iff
$c_u > c_o \cdot \alpha^*$, i.e., the participation value
exceeds the silence-channel cost.

In the classical newsvendor, the optimal order quantity is
where $F(Q^*) = c_u/(c_u + c_o)$. Here, because $V$ is
\emph{linear} in $\alpha$ (piecewise), the ``demand
distribution'' is degenerate and the optimum is at a
corner: either $\alpha = 0$ or $\alpha = \alpha^*$.

\paragraph{Extension: convex political cost.}
If the FDA faces a convex cost of deviating from ex post
optimality---e.g., $C(\alpha) = \gamma \alpha^2$---the
problem admits an interior solution:
\[
  \alpha^{**}
  = \arg\max_\alpha \; V(\alpha) - \gamma \alpha^2.
\]
With quadratic cost, the first-order condition gives a
closed-form solution analogous to the newsvendor critical
ratio, and the FDA's optimal $\alpha$ varies smoothly
with $\gamma$.


% ============================================================
\section{Monte Carlo Simulation}
\label{sec:monte_carlo}
% ============================================================

Monte Carlo methods (Session~4) are used for three purposes.

\subsection{Welfare Estimation}

For a given mechanism $\mathcal{M}$ and parameter
configuration:

\begin{algorithmic}[1]
\FOR{$k = 1, \ldots, N$}
  \STATE Draw drug quality:
    $\omega_k \sim \text{Bernoulli}(p)$
  \STATE Draw firm type:
    $\theta_k \sim \text{Bernoulli}(t)$
  \STATE \textbf{Participation:}
    compute $EU^A_{\theta_k}$; if $< c$, record
    FDA payoff $= \mathbb{1}[\omega_k = B]$ and skip
  \STATE Draw trial outcome $j_k$ from experiment
    $\pi_{\theta_k}$
  \STATE \textbf{Disclosure:} if $j_k = g$, disclose;
    if $j_k = b$, submit without evidence
  \STATE \textbf{FDA decision:} approve with prob
    $A_g^{\theta_k}$ (if disclosed) or $\alpha$
    (if no evidence)
  \STATE Record: FDA payoff $= \mathbb{1}[$decision
    matches $\omega_k]$
\ENDFOR
\STATE $\widehat{V} = N^{-1} \sum_k \text{payoff}_k$
\end{algorithmic}

Since the model has finite types and binary outcomes, the
exact expected payoff can be computed analytically (8 cells:
quality $\times$ type $\times$ outcome). Monte Carlo is
primarily useful for:

\begin{itemize}
  \item \textbf{Dynamic extension} (Section~\ref{sec:mdp}):
    where analytical solutions are unavailable.
  \item \textbf{Tail risk estimation}: probability of
    $\geq k$ bad approvals in $T$ periods.
  \item \textbf{Sensitivity analysis}: mapping $V$ as a
    function of $(p, t, c, a_1, b_1, a_2, b_2)$ when
    exact optimization is too slow for the full grid.
\end{itemize}


% ============================================================
\section{Parallel Parameter Sweep}
\label{sec:parallel}
% ============================================================

The welfare comparison (Corollary~3 in the theory paper)
requires solving the optimization across a parameter grid
(Session~5).

\subsection{Parameter Space}

\begin{itemize}
  \item $p \in \{0.05, 0.10, \ldots, 0.45\}$
    \hfill (9 values)
  \item $t \in \{0.1, 0.2, \ldots, 0.9\}$
    \hfill (9 values)
  \item $c \in \{0.1, 0.15, \ldots, 0.9\}$
    \hfill (17 values)
  \item $(a_1, b_1, a_2, b_2)$: grid over feasible set
    parameters
\end{itemize}

Total grid points: $9 \times 9 \times 17 = 1377$ (for
fixed feasible sets). Each requires solving 4 continuous
NLPs (one per participation case).

\subsection{Implementation}

Grid points are \textbf{embarrassingly parallel}: no data
dependencies. Use Python \texttt{multiprocessing.Pool}:

\begin{enumerate}
  \item \textbf{Worker function}: takes
    $(p, t, c, a_1, b_1, a_2, b_2)$, solves the
    optimization for both my model and Min's model,
    returns $(V_{\text{mine}}, V_{\text{Min}},
    \alpha^*, \pi_1^*, \pi_2^*)$.

  \item \textbf{Distribute}: \texttt{pool.map(worker,
    grid\_points)} across $C$ CPU cores.

  \item \textbf{Aggregate}: collect results into a
    DataFrame. Compute welfare cost
    $\Delta V = V_{\text{Min}} - V_{\text{mine}}$
    at each grid point.
\end{enumerate}

\subsection{Outputs}

\begin{itemize}
  \item \textbf{Heatmap}: $\Delta V(p, t)$ for fixed
    $c$ and feasible sets. Identifies parameter regions
    where cherry-picking is most costly.
  \item \textbf{Threshold surface}: $c^*(p, t)$ as a
    function of prior and type distribution.
  \item \textbf{Regime map}: which participation regime
    (both, only type~1, neither) is optimal at each
    grid point.
\end{itemize}

Speedup: $C = 8$ cores, ideal $8\times$. With
$T_1 \approx 0.1$s per grid point, total wall time
$\approx 1377 \times 0.1 / 8 \approx 17$s. Connects to
Amdahl's law (Session~5).


% ============================================================
\section{MDP Extension: Dynamic Regulation}
\label{sec:mdp}
% ============================================================

Embed the static model in a repeated setting (Sessions~2--3).

\paragraph{State:}
$s_t = (n_t^G, n_t^B)$: FDA's running count of
correct/incorrect approvals (reputation).

\paragraph{Action:}
$a_t = \alpha_t \in [0, 1]$: the commitment level for
period $t$.

\paragraph{Transition:}
Given $\alpha_t$, the static mechanism plays out.
The FDA observes the drug's true quality ex post with
some probability (post-market surveillance), updating
$(n_t^G, n_t^B)$.

\paragraph{Reward:}
\[
  r_t = EU^P(\alpha_t)
  - \gamma \cdot \mathbb{1}[\text{bad drug approved}],
\]
where $\gamma > 0$ captures reputational cost.

\paragraph{Bellman equation:}
\begin{equation}\label{eq:bellman}
  W(s) = \max_{\alpha \in [0,1]}
  \left\{
    r(s, \alpha)
    + \beta \sum_{s'} P(s' \mid s, \alpha) W(s')
  \right\}.
\end{equation}

Solve by \textbf{value iteration} on discretized
state/action spaces. The transition probabilities
$P(s'|s,\alpha)$ come from the static model's equilibrium
(which participation regime, which experiments, which
approval events occur).

No closed-form solution is available because the
static equilibrium depends on $\alpha$, which depends on
the state through the reputation cost.


% ============================================================
\section{Summary of Course Connections}
\label{sec:summary}
% ============================================================

\begin{center}
\renewcommand{\arraystretch}{1.3}
\begin{tabular}{l|c|l}
\textbf{Course Topic} & \textbf{Session}
  & \textbf{Application in This Paper} \\ \hline
Newsvendor / EOQ & 1
  & Critical ratio for $\alpha^*$ decision
    (Section~\ref{sec:newsvendor}) \\
MDP (analytical) & 2
  & Bellman equation for dynamic regulation
    (Section~\ref{sec:mdp}) \\
MDP (numerical) & 3
  & Value iteration for discretized MDP \\
Monte Carlo & 4
  & Welfare estimation and tail risk
    (Section~\ref{sec:monte_carlo}) \\
Parallel methods & 5
  & Parameter sweep across grid
    (Section~\ref{sec:parallel}) \\
\end{tabular}
\end{center}

Multi-armed bandits (Session~6) and deep RL (Session~7):
if the FDA does not know the type distribution $t$ and
must learn through repeated interactions, the problem
becomes a bandit/RL problem. This is a possible extension
but is not developed here.


\end{document}
